% TODO proof-reading
\subsubsection{二次元配列}

Javaの二次元配列は、別の一次元配列への\IT{参照}の一次元配列に過ぎません。


二次元配列を作成してみましょう：

\begin{lstlisting}[style=customjava]
	public static void main(String[] args)
	{
		int[][] a = new int[5][10];
		a[1][2]=3;
	}
\end{lstlisting}

\begin{lstlisting}
  public static void main(java.lang.String[]);
    flags: ACC_PUBLIC, ACC_STATIC
    Code:
      stack=3, locals=2, args_size=1
         0: iconst_5      
         1: bipush        10
         3: multianewarray #2,  2      // class "[[I"
         7: astore_1      
         8: aload_1       
         9: iconst_1      
        10: aaload        
        11: iconst_2      
        12: iconst_3      
        13: iastore       
        14: return        
\end{lstlisting}

\TT{multianewarray}命令を使用して作成されます：オブジェクトの型と次元がオペランドとして渡されます。

配列のサイズ（10*5）はスタックに残されます（\TT{iconst\_5}と\TT{bipush}命令を使用）。


行\#1への\IT{参照}はオフセット10でロードされます（\TT{iconst\_1}と\TT{aaload}）。

列はオフセット11の\TT{iconst\_2}を使用して選択されます。

書き込まれる値はオフセット12で設定されます。

13の\TT{iastore}は配列の要素を書き込みます。


要素はどのようにアクセスされるのでしょうか？

\begin{lstlisting}[style=customjava]
	public static int get12 (int[][] in)
	{
		return in[1][2];
	}
\end{lstlisting}

\begin{lstlisting}
  public static int get12(int[][]);
    flags: ACC_PUBLIC, ACC_STATIC
    Code:
      stack=2, locals=1, args_size=1
         0: aload_0       
         1: iconst_1      
         2: aaload        
         3: iconst_2      
         4: iaload        
         5: ireturn       
\end{lstlisting}

配列の行への\IT{参照}はオフセット2でロードされ、列はオフセット3で設定され、その後\TT{iaload}が配列の要素をロードします。